\documentclass{article}

\usepackage{amsmath,amssymb,amsthm}
\usepackage{fullpage}
\usepackage{mathtools}

\newcommand{\RR}{\mathbb{R}}
\newcommand{\CC}{\mathbb{C}}
\newcommand{\QQ}{\mathbb{Q}}
\newcommand{\NN}{\mathbb{N}}
\newcommand{\ZZ}{\mathbb{Z}}
\newcommand{\PP}{\mathcal{P}}
\DeclareMathOperator{\GL}{GL}
\DeclareMathOperator{\SL}{SL}
\DeclareMathOperator{\cis}{cis}

\theoremstyle{definition}
\newtheorem*{solution}{Solution}

\title{Math 430 -- Problem Set 6 Solutions}
\date{Due April 18, 2016}

\begin{document}
\maketitle


\begin{enumerate}
\item[16.27.] Let $R$ be a commutative ring.  An element $a$ in $R$ is nilpotent if $a^n = 0$ for some positive integer $n$.  Show that the set of all nilpotent elements forms an ideal in $R$.
\begin{solution} 
Let $N \subseteq R$ be the set of nilpotent elements, and suppose $a^m = 0$ and $b^n = 0$.  Then $(-a)^m = 0$, so $N$ is closed under negation.  Moreover, $(a + b)^{m+n} = \sum_{i=0}^{m+n} \binom{m+n}{i} a^i b^{m+n-i} = 0$ since either $i \ge m$ and $a^i = 0$ or $m+n-i \ge n$ and $b^{m+n-i} = 0$.  Thus $N$ is closed under addition.  Finally, if $x \in R$ then $(ax)^m = a^mx^m = 0$, so $N$ is an ideal.
\end{solution}
\item[16.40.] Let $R$ be a ring and $I$ and $J$ be ideals in $R$ such that $I + J = R$.
\begin{enumerate}
\item Show that for any $r$ and $s$ in $R$, the system of equations
\begin{align*}
x &\equiv r \pmod{I} \\
x &\equiv s \pmod{J}
\end{align*}
has a solution.
\begin{solution}
Since $I + J = R$, we may find $i \in I$ and $j \in J$ with $i + j = 1$.  Setting $x = si + rj$, we have
\begin{align*}
x &\equiv rj \equiv r \pmod{I} \\
x &\equiv si \equiv s \pmod{J}.
\end{align*}
\end{solution}
\item In addition, prove that any two solutions of the system are congruent modulo $I \cap J$.
\begin{solution}
If $x$ and $y$ are solutions, then $x - y \equiv 0 \pmod{I}$ and $x - y \equiv 0 \pmod{J}$, so $x - y \in I \cap J$.
\end{solution}
\item Let $I$ and $J$ be ideals in a ring $R$ such that $I + J = R$. Show that there exists a ring isomorphism
\[
R / (I \cap J) \cong R/I \times R/J.
\]
\begin{solution}
Let $\phi : R \to R/I \times R/J$ be defined by $\phi(x) = (x + I, x + J)$.  By part (a), $\phi$ is surjective, and by part (b) it has kernel $I \cap J$.  So by the First Isomorphism Theorem, it induces the desired isomorphism.
\end{solution}
\end{enumerate}
\item[17.2. (b)] Compute $(5x^2 + 3x - 4)(4x^2 - x + 9)$ in $\ZZ_{12}[x]$.
\begin{solution}
$8x^4 + 7x^3 + 2x^2 + 7x$.
\end{solution}
\item[17.3. (b)] Let $a(x) = 6x^4 - 2x^3 + x^2 - 3x + 1$ and $b(x) = x^2 + x - 2$ in $\ZZ_7[x]$.  Use the division algorithm to find $q(x)$ and $r(x)$ so that $a(x) = q(x)b(x) + r(x)$ with $\deg r(x) < \deg b(x)$.
\begin{solution}
We find that $q(x) = 6x^2 + 6x$ and $r(x) = 2x+1$.
\end{solution}
\item[17.4. (c)] Find the greatest common divisor $d(x)$ of $p(x) = x^3 + x^2 - 4x + 4$ and $q(x) = x^3 + 3x - 2$ in $\ZZ_5[x]$ and polynomials $a(x)$ and $b(x)$ such that $a(x)p(x) + b(x)q(x) = d(x).$
\begin{solution}
\begin{align*}
x^3 + x^2 - 4x + 4 &= 1(x^3 + 3x - 2) + (x^2 + 3x + 1) \\
x^3 + 3x - 2 &= (x+2)(x^2 + 3x + 1) + (x+1) \\
x^2 + 3x + 1 &= (x+2)(x+1) + 4 \\
4 &= (x^2 + 3x + 1) - (x+2)(x+1) \\
&= (x^2 + 3x + 1) - (x+2)((x^3 +3x - 2) - (x+2)(x^2 + 3x+1)) \\
&= (x^2 + 4x)(x^2 + 3x + 1) - (x+2)(x^3 + 3x - 2) \\
&= (x^2 + 4x)((x^3 + x^2 -4x + 4) - (x^3 + 3x - 2)) - (x+2)(x^3 + 3x - 2) \\
&= (x^2 + 4x)(x^3 + x^2 -4x + 4) - (x^2 + 2)(x^3 + 3x - 2)
\end{align*}
Negating this last equation, we get
\[
1 = (4x^2 + x)(x^3 + x^2 -4x + 4) + (x^2 + 2)(x^3 + 3x - 2).
\]
\end{solution}
\item[17.7.] Find a unit $p(x)$ in $\ZZ_4[x]$ such that $\deg p(x) > 1$.
\begin{solution}
Since $(2x^2+1)(2x^2+1) = 1$, we may take $p(x) = 2x^2 + 1$.
\end{solution}
\item[17.9.] Find all of the irreducible polynomials of degrees $2$ and $3$ in $\ZZ_2[x]$.
\begin{solution}
A polynomial of degree $2$ or $3$ is irreducible when it has no roots.  The only two possible roots in $\ZZ_2$ are $0$ and $1$, so $f(x) = x^2 + ax + b$ is irreducible when $f(0) = b = 1$ and $f(1) = 1 + a + b = 1$, so $a = b =1$.

Similarly, $f(x) = x^3 + ax^2 + bx + c$ is irreducible when $g(0) = c = 1$ and $g(1) = 1 + a + b + c = 1$, yielding $a = 1$ and $b= 0$ or $a = 0$ and $b = 1$.  Thus the irreducible polynomials are
\[
x^2 + x + 1, \\
x^3 + x + 1, \\
x^3 + x^2 + 1.
\]
\end{solution}
\item[17.10.] Give two different factorizations of $x^2 + x + 8$ in $\ZZ_{10}[x]$.
\begin{solution}
We first find the roots by trial and error: $x = 1, -2, 3, -4$.  Pairing these up so that they have product $-2$ and sum $1$, we get the factorizations
\begin{align*}
x^2 + x + 8 &= (x - 1)(x + 2) \\
&= (x - 3)(x + 4).
\end{align*}
\end{solution}
\item[17.18.] Let $p(x) = a_nx^n + x_{n-1}x^{n-1} + \dots + a_0 \in \ZZ[x]$, with $a_n \ne 0$.  Prove that if $p(r/s) = 0$ with $\gcd(r, s) = 1$, then $r \mid a_0$ and $s \mid a_n$.
\begin{solution}
Substituting $r/s$ into $p(x)$ and multiplying by $s^n$ we get
\begin{align*}
0 &= a_nr^n + a_{n-1}r^{n-1}s + \cdots + a_1 r s^{n-1} + a_0s^n \\
a_nr^n &= s(-a_{n-1}r^{n-1} - \cdots - a_1 r s^{n-2} - a_0 s^{n-1}) \\
a_0s^n &= r(-a_nr^{n-1} - a_{n-1}r^{n-2} - \cdots - a_1 s^{n-1}).
\end{align*}
Thus $s$ divides $a_nr^n$ and $r$ divides $a_0s^n$.  Since $r$ and $s$ are relatively prime, $s$ divides $a_n$ and $r$ divides $a_0$.
\end{solution}
\item[17.20.] Let $\Phi_n(x) = \frac{x^n-1}{x-1} = x^{n-1} + x^{n-2} + \dots + x + 1.$ Show that $\Phi_p(x)$ is irreducible over $\QQ$ for any prime $p$.
\begin{solution}
We make the substitution $t = x-1$, yielding
\begin{align*}
\Phi_p(x) &= \frac{(t+1)^p - 1}{t} \\
&= \sum_{i = 1}^p \binom{p}{i} t^{i-1}.
\end{align*}
Since $\binom{p}{i}$ is divisible by $p$ for $0 < i < p$ and $\binom{p}{1} = p$ is not divisible by $p^2$ and $\binom{p}{p} = 1$ is not divisible by $p$, this polynomial satisfies the Eisenstein criterion and is thus irreducible.  Thus $\Phi_p(x)$ is irreducible as well.
\end{solution}
\item[17.21.] If $F$ is a field, show that there are infinitely many irreducible polynomials in $F[x]$.
\begin{solution}
Euclid's proof for the infinitude of primes in $\ZZ$ applies in essentially the same way here.  Suppose that there were finitely many irreducible polynomials $p_1, \dots, p_k$.  Let $p = 1 + \prod_{i=1}^k p_i$.  Since $p$ has remainder $1$ when divided by each $p_i$, it is not a multiple of any of them.  But it must be divisible by some irreducible since $F[x]$ is Noetherian.  Thus there are infinitely many irreducible polynomials.
\end{solution}
\item[17.24.] Show that $x^p - x$ has $p$ distinct zeros in $\ZZ_p$, for any prime $p$.  Conclude that
\[
x^p - x = x(x-1)(x-2) \cdots (x - (p-1)).
\]
\begin{solution}
By Fermat's little theorem, $a^p \equiv a \pmod{p}$ for all $a \in \ZZ_p$ and thus $x^p - x$ is divisible by $(x-a)$ for all $a \in \ZZ_p$.  By additivity of degree and the equality of leading coefficients, we get the desired equation.
\end{solution}
\end{enumerate}

\end{document}